% Generated draft: requires mathematical and source review.
\documentclass[11pt]{article}
\usepackage[margin=25mm]{geometry}
\usepackage{fontspec}
\setmainfont{Latin Modern Roman}
\usepackage{amsmath,amssymb,mathtools,unicode-math}
\setmathfont{Latin Modern Math}
\usepackage{enumitem,needspace,xurl,hyperref}
\setlength{\parindent}{0pt}
\setlength{\parskip}{5pt}
\setlength{\emergencystretch}{3em}
\newcommand{\N}{\mathbb{N}}
\newcommand{\Z}{\mathbb{Z}}
\newcommand{\Q}{\mathbb{Q}}
\newcommand{\R}{\mathbb{R}}
\newcommand{\C}{\mathbb{C}}
\newcommand{\F}{\mathbb{F}}
\newcommand{\A}{\mathbb{A}}
\newcommand{\PP}{\mathbb{P}}
\newcommand{\E}{\mathbb{E}}
\newcommand{\eps}{\varepsilon}
\newcommand{\dd}{\,\mathrm{d}}
\begin{document}
\section*{UnsolvedMath ID: 700574}
\textbf{Source problem code:} catalogue entry
\subsection{Definitions and mathematical statement}
Let \(\Omega\) be a countably infinite set, and let \(\mathrm{Sym}(\Omega)\) be its group of all permutations. For a cardinal \(\kappa\in\{2,3,\ldots\}\cup\{\aleph_0\}\), define the set of \(\kappa\)-sections by
\[
\mathcal S_\kappa(\Omega)=\left\{\mathcal P\subseteq\mathcal P(\Omega): |\mathcal P|=\kappa,\ \bigcup_{A\in\mathcal P}A=\Omega,\ A\cap B=\varnothing\text{ for distinct }A,B\in\mathcal P,\ |A|=\aleph_0\text{ for every }A\in\mathcal P\right\}.
\]
Thus the blocks of a section are not labeled. Put \(I_k=\{1,\ldots,k\}\) for finite \(k\geq2\), and \(I_{\aleph_0}=\N\). An ordered \(\kappa\)-section is an indexed family \((A_i)_{i\in I_\kappa}\) whose distinct blocks form a member of \(\mathcal S_\kappa(\Omega)\).

A subgroup \(G\leq\mathrm{Sym}(\Omega)\) acts on sections by \(g\mathcal P=\{g[A]:A\in\mathcal P\}\), and on ordered sections by \(g(A_i)_{i\in I_\kappa}=(g[A_i])_{i\in I_\kappa}\). Transitivity on sections means that for every \(\mathcal P,\mathcal Q\in\mathcal S_\kappa(\Omega)\) some \(g\in G\) satisfies \(g\mathcal P=\mathcal Q\). Transitivity on ordered sections requires, for every two such indexed families, a single \(g\in G\) satisfying \(g[A_i]=B_i\) for every index \(i\). Transitivity on \(\Omega\) means that for all \(x,y\in\Omega\) some \(g\in G\) satisfies \(g(x)=y\).

The supplied record contains the following questions, retaining its part labels.
\begin{enumerate}
\item[(b)] For a specified finite integer \(k\geq2\), does there exist \(G\leq\mathrm{Sym}(\Omega)\) that is transitive on \(\mathcal S_k(\Omega)\) but not transitive on \(\mathcal S_{k+1}(\Omega)\)? Here failure of transitivity means that there exist \(\mathcal P,\mathcal Q\in\mathcal S_{k+1}(\Omega)\) such that \(g\mathcal P\ne\mathcal Q\) for every \(g\in G\). This is the finite-parameter reading of the source: the source does not explicitly specify the range or outer quantifier for \(k\), or require \(G\) to be transitive on \(\Omega\) in this part. Consequently, whether existence is requested for every finite \(k\), or for some \(k\), remains unresolved. If \(k=\aleph_0\) were included and \(k+1\) meant cardinal addition, the two section cardinalities would coincide.
\item[(c)] Does there exist \(G\leq\mathrm{Sym}(\Omega)\) such that \(G\) is transitive on \(\Omega\), transitive on \(\mathcal S_{\aleph_0}(\Omega)\), and \(G\ne\mathrm{Sym}(\Omega)\)?
\item[(a), archived] For a specified \(\kappa\in\{2,3,\ldots\}\cup\{\aleph_0\}\), does there exist a subgroup \(G\leq\mathrm{Sym}(\Omega)\) transitive on \(\Omega\) and on \(\mathcal S_\kappa(\Omega)\), but not transitive on ordered \(\kappa\)-sections? The inclusion of \(\aleph_0\) here follows the explicit range in the supplied archive commentary; the question text itself only writes \(k\geq2\).
\end{enumerate}

The supplied archive attributes an affirmative answer to (a) to P. M. Neumann and gives the construction
\[
G_U=\{g\in\mathrm{Sym}(\Omega):\mathrm{Fix}(g)\in U\},\qquad
\mathrm{Fix}(g)=\{x\in\Omega:g(x)=x\},
\]
where \(U\) is a nonprincipal ultrafilter on \(\Omega\). Here an ultrafilter is a family \(U\subseteq\mathcal P(\Omega)\) containing \(\Omega\), excluding \(\varnothing\), closed under finite intersections and supersets, and containing exactly one of \(A\) and \(\Omega\setminus A\) for every \(A\subseteq\Omega\); nonprincipal means that it contains no finite subset. The archive asserts that this single group has the required section and ordered-section behavior for every \(2\leq\kappa\leq\aleph_0\). This is a supplied historical assertion, not a proof or an independently verified result in this draft.

The supplied main-section commentary labels (c) solved affirmatively, describes a 2026 AI solution by Achyuth Jayadevan as confirmed by S. M. Corson, and also reports a 2025 result of S. M. Corson and S. Shelah establishing consistency of an affirmative answer with ZFC. These are distinct supplied claims: consistency with ZFC is not itself an assertion of a proof in ZFC. The registry labels the overall record partially solved; no independent status assessment is made here.
\subsection{Short English statement}
For permutations of a countably infinite set, ask whether transitivity on partitions into \(k\) infinite blocks can hold while transitivity on partitions into \(k+1\) infinite blocks fails. The record also asks for a proper, point-transitive subgroup transitive on countably infinite partitions, and archives the question of transitivity on unlabeled partitions without transitivity on labeled partitions. Parts (a) and (c) have supplied affirmative solution reports; the parameter quantifier in (b) remains unclear.
\subsection{Sources}
\begin{itemize}
\item Local UnsolvedMath registry, permanent ID \(700574\), exact source problem code \texttt{catalogue entry}, title ``Kourovka Notebook Problem 9.41.'' The supplied complete source record attributes the problem to P. M. Neumann, proposed in 1984.
\item Supplied primary-source attribution: \emph{The Kourovka Notebook (Unsolved Problems in Group Theory)}, 9th issue (1984), problem 9.41, updated in the 21st issue (2026), \texttt{arXiv:1401.0300v48}. The supplied collection description names E. I. Khukhro and V. D. Mazurov as editors. These source contents have not been independently consulted for this draft.
\item Supplied Archive of Solved Problems entry for part (a): P. M. Neumann, \emph{Letter of October, 5, 1989}. The entry reports an affirmative answer using a nonprincipal ultrafilter and asserts the stated behavior for every \(2\leq\kappa\leq\aleph_0\).
\item Supplied main-section solution report for part (c): Achyuth Jayadevan, \emph{Preprint}, 2026, \url{https://zenodo.org/records/22820556}. The commentary describes this as an AI solution confirmed by S. M. Corson; that description and confirmation are supplied attribution, not independently verified here.
\item Supplied consistency report for part (c): S. M. Corson and S. Shelah, \emph{Preprint}, 2025, \url{https://arxiv.org/abs/2503.12997}. The commentary reports consistency of an affirmative answer with the ZFC axioms of set theory.
\item Supplied registry status review dated \texttt{2026-10-16}: canonical status \texttt{partially\_solved}, evidence field ``published,'' with a note referring to archived solved parts and remaining main-section parts. This is supplied dated curation, not an independently verified mathematical or literature assessment.
\end{itemize}
\end{document}
